% ============================================================
%  P.T. Tsilingiris, "Thermophysical and transport properties of
%  humid air at temperature range between 0 and 100 degC"
%  Energy Conversion and Management 49 (2008) 1098-1110
%  doi:10.1016/j.enconman.2007.09.015
%
%  LaTeX transcription: structure, nomenclature, all numbered
%  equations, all numerical constants, all data tables, and the
%  reference list.  Body prose is marked with placeholder
%  comments -- see the note accompanying this file.
%
%  NOTE ON TWO DEFECTS IN THE ORIGINAL, both flagged inline below:
%   (a) Table 4 prints SK2 = -1.788037411E-2.  That value is
%       physically impossible (it drives k_m to -178 W/m K at
%       100 degC).  Numerical check against the paper's own
%       Eq. (29) confirms the intended value is
%       SK2 = -1.788037411E-7, which reproduces Eq. (29) to
%       ~1% and matches Fig. 3 (0.0240 at 0 degC, maximum near
%       60 degC).  See the comment at Table 4.
%   (b) Eq. (39) is captioned "in W/m K x 10^-3", but its
%       coefficients already return W/m K directly
%       (k_a(273 K) = 0.02403).  Eq. (42) for water vapor DOES
%       need the 10^-3 factor.  See the comment at Eq. (39).
%
%  Compile: pdflatex tsilingiris2008.tex   (twice)
% ============================================================

\documentclass[11pt]{article}

\usepackage[T1]{fontenc}
\usepackage[utf8]{inputenc}
\usepackage{mathptmx}
\usepackage{amsmath,amssymb}
\usepackage{booktabs}
\usepackage{graphicx}
\usepackage{caption}
\usepackage{array}
\usepackage{longtable}
\usepackage{pdflscape}
\usepackage[margin=1in]{geometry}
\usepackage{fancyhdr}
\usepackage[hidelinks]{hyperref}

\pagestyle{fancy}
\fancyhf{}
\fancyhead[C]{\itshape P.T. Tsilingiris / Energy Conversion and Management 49 (2008) 1098--1110}
\fancyfoot[C]{\thepage}
\renewcommand{\headrulewidth}{0pt}

% --- notation shortcuts -------------------------------------
\newcommand{\degC}{$^{\circ}$C}
\newcommand{\mdegC}{^{\circ}\mathrm{C}}
\newcommand{\Psv}{P_\mathrm{sv}}
\newcommand{\RH}{\mathrm{RH}}
\newcommand{\fPT}{f(P,T)}
% The recurring group  f(P,T).RH.(Psv/P0)  appears in Eqs.
% (21), (29), (32) and (34); \grp keeps them readable.
\newcommand{\grp}{\fPT\cdot\RH\cdot\left(\dfrac{\Psv}{P_0}\right)}
\newcommand{\frcaption}[1]{\begin{center}\itshape #1\end{center}}

\begin{document}

\title{Thermophysical and transport properties of humid air at
       temperature range between 0 and 100 \degC}
\author{P.T. Tsilingiris\thanks{Tel.: +30 210 9355289.
        \textit{E-mail address:} ptsiling@teiath.gr} \\[4pt]
        \small\itshape Department of Energy Engineering, Technological
        Education Institution (TEI) of Athens, \\
        \small\itshape A. Spyridonos Street, GR 122 10 Egaleo, Athens, Greece}
\date{\small Received 4 February 2007; accepted 29 September 2007 \\
      Available online 28 November 2007 \\[6pt]
      \normalsize Energy Conversion and Management \textbf{49} (2008) 1098--1110}
\maketitle

\begin{abstract}
% [Abstract text -- paste from the original, p.\ 1098]
\end{abstract}

\noindent\textit{Keywords:} Thermophysical properties; Density;
Viscosity; Specific heat capacity; Thermal conductivity; Thermal
diffusivity; Prandtl number

% ============================================================
\section*{Nomenclature}
% ============================================================

\subsection*{Symbols}
\begin{tabbing}
\hspace{3.6cm}\= \kill
$a$, $b$, $c$ \> numerical constants \\
$A$, $B$ \> virial coefficients \\
$A_0$, $A_1$, $A_2$, $A_3$ \> numerical constants \\
$B_0$, $B_1$, $B_2$, $B_3$ \> numerical constants \\
$c_p$ \> specific heat capacity (J/kg\,K) \\
$C_1$, $C_2$, $C_3$ \> numerical constants \\
$\mathrm{CA}_0$, $\mathrm{CA}_1$, \ldots, $\mathrm{CA}_4$ \> numerical constants \\
$\mathrm{CV}_0$, $\mathrm{CV}_1$, $\mathrm{CV}_2$ \> numerical constants \\
COD \> coefficient of determination \\
$E_0$, $E_1$, $E_2$, $E_3$, $E_4$ \> numerical constants \\
$f$ \> enhancement factor \\
$k$ \> thermal conductivity (W/m\,K) \\
$K_1$, $K_2$, $K_3$ \> numerical constants \\
$\mathrm{KA}_0$, $\mathrm{KA}_1$, \ldots, $\mathrm{KA}_5$ \> numerical constants \\
$\mathrm{KV}_0$, $\mathrm{KV}_1$, $\mathrm{KV}_2$ \> numerical constants \\
$M$ \> molar mass (kg/kmol) \\
$\mathrm{MA}_0$, $\mathrm{MA}_1$, \ldots, $\mathrm{MA}_4$ \> numerical constants \\
$\mathrm{MV}_0$, $\mathrm{MV}_1$ \> numerical constants \\
$n$ \> mole number \\
$P$ \> pressure (Pa) \\
$Pr$ \> Prandtl number \\
$R$ \> universal gas constant (8.31441 J/mol\,K) \\
$\RH$ \> relative humidity \\
$\mathrm{SD}_0$, $\mathrm{SD}_1$, \ldots, $\mathrm{SD}_3$ \> numerical constants \\
$\mathrm{SV}_0$, $\mathrm{SV}_1$, \ldots, $\mathrm{SV}_4$ \> numerical constants \\
$\mathrm{SC}_0$, $\mathrm{SC}_1$, \ldots, $\mathrm{SC}_6$ \> numerical constants \\
$\mathrm{SA}_0$, $\mathrm{SA}_1$, \ldots, $\mathrm{SA}_5$ \> numerical constants \\
$\mathrm{SK}_0$, $\mathrm{SK}_1$, \ldots, $\mathrm{SK}_5$ \> numerical constants \\
$\mathrm{SP}_0$, $\mathrm{SP}_1$, \ldots, $\mathrm{SP}_4$ \> numerical constants \\
$t$ \> temperature (\degC) \\
$T$ \> absolute temperature (K) \\
$z$ \> compressibility factor
\end{tabbing}

\subsection*{Greek letters}
\begin{tabbing}
\hspace{3.6cm}\= \kill
$\alpha$ \> thermal diffusivity (m$^2$/s) \\
$\Delta$ \> difference \\
$\epsilon$ \> dimensionless multiplier \\
$\Theta$ \> interaction parameter \\
$\mu$ \> viscosity (Ns/m$^2$) \\
$\xi_1$, $\xi_2$ \> dimensionless multipliers \\
$\rho$ \> density (kg/m$^3$) \\
$\Phi$ \> interaction parameter
\end{tabbing}

\subsection*{Subscripts}
\begin{tabbing}
\hspace{1.6cm}\= \kill
a  \> air \\
i  \> ideal \\
m  \> mixture \\
sv \> saturated vapor \\
tr \> monoatomic value \\
v  \> vapor \\
0  \> total
\end{tabbing}

% ============================================================
\section{Introduction}
% ============================================================
% [Introduction prose -- pp.\ 1098-1099.  Key points: prior work
%  targets either low temperatures (Giacomo [1], Davies [2],
%  Zuckerwar and Meredith [3], Rasmussen [4], Hyland and Wexler
%  [5]) or 100-200 degC (Melling et al. [6]); the 0-100 degC band
%  matters for drying and water distillation; apart from Nelson's
%  brief survey [7], correlations for humid-air transport
%  properties in this range are lacking.  Substituting dry-air for
%  moist-air properties can cause appreciable error near
%  saturation because M_v < M_a.]

% ============================================================
\section{The evaluation of thermophysical and transport properties
         of humid air}
% ============================================================

Humid air is regarded as a binary mixture of dry air and water
vapor. The molar fraction of water vapor is defined as the ratio of
water vapor moles to the total number of moles of the mixture as
\begin{equation}
x_\mathrm{v} = \frac{n_\mathrm{v}}{n_\mathrm{m}}
             = \frac{n_\mathrm{v}}{n_\mathrm{a}+n_\mathrm{v}}
             = \frac{P_\mathrm{v}}{P_0}
\end{equation}
The relative humidity is defined as
\begin{equation}
\RH = \frac{n_\mathrm{v}}{n_\mathrm{sv}}
    = \frac{x_\mathrm{v}}{x_\mathrm{sv}}
    = \frac{P_\mathrm{v}}{\Psv}
\end{equation}
from which
\begin{equation}
x_\mathrm{v} = x_\mathrm{sv}\cdot\RH
\end{equation}

% [Prose -- p.\ 1100, on the enhancement factor [1,8].]

\noindent This is introduced to correct the molar fraction of
saturated vapor pressure,
\begin{equation}
x_\mathrm{sv} = \fPT\cdot\frac{\Psv}{P_0}
\end{equation}
The molar fraction of water vapor is then calculated from Eqs.~(3)
and (4) as a function of the total atmospheric pressure $P_0$ and the
saturated vapor pressure $\Psv$ at a specific temperature by the
following expression,
\begin{equation}
x_\mathrm{v} = \fPT\cdot\RH\cdot\frac{\Psv}{P_0}
\end{equation}

\noindent Its calculation was performed according to Hardy [11] by the
following simplified fitting expression recommended by
Greenspan [12]
\begin{equation}
\fPT = \exp\left[\xi_1\cdot\left(1-\frac{\Psv}{P_0}\right)
       + \xi_2\cdot\left(\frac{P_0}{\Psv}-1\right)\right]
\end{equation}
with
\begin{equation}
\xi_1 = \sum_{i=0}^{3} A_i\cdot T^{i}
\end{equation}
\begin{equation}
\xi_2 = \exp\left[\sum_{i=0}^{3} B_i\cdot T^{i}\right]
\end{equation}

\noindent The numerical values of the constants in Eqs.~(7) and (8)
corresponding to the temperature range between 0 and 100 \degC\ are
$A_0 = 3.53624\times10^{-4}$,
$A_1 = 2.93228\times10^{-5}$,
$A_2 = 2.61474\times10^{-7}$,
$A_3 = 8.57538\times10^{-9}$,
$B_0 = -1.07588\times10^{1}$,
$B_1 = 6.32529\times10^{-2}$,
$B_2 = -2.53591\times10^{-4}$ and
$B_3 = 6.33784\times10^{-7}$.

% [Prose -- p.\ 1100.  The f = 1 assumption leads to less than 0.5%
%  maximum error at temperatures around 75 degC.]

\noindent Values of $\Psv$ were derived by the following fourth degree
polynomial, fitted to saturation vapor pressure data between 0 and
100 \degC\ from the thermodynamic properties of water [14],
\begin{equation}
\Psv = E_0 + E_1\cdot t + E_2\cdot t^{2} + E_3\cdot t^{3}
       + E_4\cdot t^{4}
\end{equation}
where $\Psv$ is in kPa, for the following values of numerical
constants
\linebreak
$E_0 = 0.7073034146$,
$E_1 = -2.703615165\times10^{-2}$,
$E_2 = 4.36088211\times10^{-3}$,
$E_3 = -4.662575642\times10^{-5}$ and
$E_4 = 1.034693708\times10^{-6}$.\hfill\null

% ------------------------------------------------------------
\subsection{Density}

\noindent The density of the binary mixture of pure water vapor and dry
air at the corresponding partial pressures and molar fractions of
$P_\mathrm{v}$, $x_\mathrm{v}$ and
$P_\mathrm{a} = P_0 - P_\mathrm{v}$,
$x_\mathrm{a} = 1 - x_\mathrm{v}$, respectively, is calculated through
the gas equation of state by the following simple mixing correlation,
\begin{equation}
\rho_\mathrm{m} = \frac{1}{z_\mathrm{m}(x_\mathrm{v},T)}
  \cdot\frac{P_0}{R\cdot T}
  \cdot\left( M_\mathrm{a}\cdot\frac{P_0-P_\mathrm{v}}{P_0}
            + M_\mathrm{v}\frac{P_\mathrm{v}}{P_0} \right)
\end{equation}
where $z_\mathrm{m}(x_\mathrm{v},T)$ is the compressibility factor for
the gas mixture. From the above expression, the density can be derived
as a function of the molar fraction of water vapor as
\begin{equation}
\rho_\mathrm{m} = \frac{1}{z_\mathrm{m}(x_\mathrm{v},T)}
  \cdot\frac{P_0}{R\cdot T}\cdot M_\mathrm{a}
  \cdot\left[1 - x_\mathrm{v}\cdot
    \left(1-\frac{M_\mathrm{v}}{M_\mathrm{a}}\right)\right]
\end{equation}
which, combined with Eq.~(5), leads to the following expression for the
density of the binary mixture,
\begin{equation}
\rho_\mathrm{m} = \frac{1}{z_\mathrm{m}(x_\mathrm{v},T)}
  \cdot\frac{P_0}{R\cdot T}\cdot M_\mathrm{a}
  \cdot\left[1 - \fPT\cdot\RH\cdot
    \left(1-\frac{M_\mathrm{v}}{M_\mathrm{a}}\right)
    \cdot\left(\frac{\Psv}{P_0}\right)\right]
\end{equation}

\noindent Melling et al.\ [6] calculated the compressibility factor of
humid air between 100 and 200 \degC\ by the following approximate
mixing expression,
\begin{equation}
z_\mathrm{m}(x_\mathrm{v},T) = 1 + x_\mathrm{v}
  \cdot\left(\frac{a+c\cdot T}{1+b\cdot T}-1\right)
\end{equation}

\noindent For the present investigation the compressibility factor was
evaluated from the virial equation of state according to the following
expression,
\begin{equation}
z_\mathrm{v} = 1 + A\cdot \Psv + B\cdot \Psv^{2}
\end{equation}
which was recommended by Hyland and Wexler [8]. The second and third
pressure series virial coefficients were calculated by
\begin{equation}
A = C_1 + C_2\cdot \mathrm{e}^{\frac{C_3}{T}}
\end{equation}
\begin{equation}
B = K_1 + K_2\,\mathrm{e}^{\frac{K_3}{T}}
\end{equation}
with
$C_1 = 0.7\times10^{-8}$ Pa$^{-1}$,
$C_2 = -0.147184\times10^{-8}$ Pa$^{-1}$,
$C_3 = 1734.29$ (K$^{-1}$),
$K_1 = 0.104\times10^{-14}$ Pa$^{-2}$,
$K_2 = -0.335297\times10^{-17}$ Pa$^{-2}$ and
$K_3 = 3645.09$ K$^{-1}$.

% [Prose -- p.\ 1101.  Fixing z_v = 1 costs about 0.38% at 50 degC
%  and under 1.5% at 100 degC; the combined effect of fixing both
%  z and f at unity gives an overall maximum density error between
%  about 0.4% and 1.5% at 0 and 100 degC respectively.]

% ---------------- Table 1 ----------------
\begin{table}[htbp]
\centering
\footnotesize
\caption*{Table 1\\ The calculated enhancement and compressibility
          factors for temperatures up to 100 \degC}
\setlength{\tabcolsep}{5pt}
\begin{tabular}{lcccccccccc}
\toprule
$T$ (\degC) & 10 & 20 & 30 & 40 & 50 & 60 & 70 & 80 & 90 & 100 \\
\midrule
$f$            & 1.0006 & 1.0010 & 1.0016 & 1.0023 & 1.0031
               & 1.0039 & 1.0046 & 1.0046 & 1.0034 & 1.0000 \\
$z_\mathrm{v}$ & 0.9991 & 0.9987 & 0.9981 & 0.9972 & 0.9961
               & 0.9947 & 0.9928 & 0.9906 & 0.9878 & 0.9844 \\
\bottomrule
\end{tabular}
\end{table}

% ------------------------------------------------------------
\subsection{Viscosity}

Reid et al.\ [15] recommended the following expression for the
viscosity of a mixture of dilute gases with $i$ components, which was
based on earlier investigations by Wilke [16],
\begin{equation}
\mu_\mathrm{m} = \sum_{i=1}^{n}
  \frac{x_i\cdot\mu_i}{\sum_{j=1}^{n} x_j \Phi_{ij}}
\end{equation}
with the interaction parameters $\Phi_{ij}$ and $\Phi_{ji}$ given by
\begin{equation}
\Phi_{ij} = \frac{\left[1
  + \left(\dfrac{\mu_i}{\mu_j}\right)^{1/2}
  \cdot\left(\dfrac{M_j}{M_i}\right)^{1/4}\right]^{2}}
  {\left[8\left(1+\dfrac{M_i}{M_j}\right)\right]^{1/2}}
\end{equation}
\begin{equation}
\Phi_{ji} = \frac{\mu_j}{\mu_i}\cdot\frac{M_i}{M_j}\cdot\Phi_{ij}
\end{equation}

\noindent As derived from Eq.~(18), $\Phi_{ii} = \Phi_{jj} = 1$, which,
combined with Eq.~(17), after its expansion, leads to the following
expression,
\begin{equation}
\mu_\mathrm{m} = \frac{(1-x_\mathrm{v})\cdot\mu_\mathrm{a}}
  {(1-x_\mathrm{v}) + x_\mathrm{v}\cdot\Phi_\mathrm{av}}
  + \frac{x_\mathrm{v}\cdot\mu_\mathrm{v}}
  {x_\mathrm{v} + (1-x_\mathrm{v})\cdot\Phi_\mathrm{va}}
\end{equation}

\noindent Taking into consideration Eq.~(5), the previous expression
becomes
\begin{equation}
\begin{split}
\mu_\mathrm{m} = {} &
 \frac{\left[1-\grp\right]\cdot\mu_\mathrm{a}}
      {\left[1-\grp\right] + \grp\,\Phi_\mathrm{av}} \\[6pt]
 & + \frac{\grp\,\mu_\mathrm{v}}
        {\grp + \left[1-\grp\right]\cdot\Phi_\mathrm{va}}
\end{split}
\end{equation}
which offers the viscosity of a humid air mixture at a specific
temperature and relative humidity as a function of dry air and water
vapor viscosities for the following values of the interaction
parameters,
\begin{equation}
\Phi_\mathrm{av} = \frac{\sqrt{2}}{4}
  \cdot\left(1+\frac{M_\mathrm{a}}{M_\mathrm{v}}\right)^{-\frac{1}{2}}
  \cdot\left[1
  + \left(\frac{\mu_\mathrm{a}}{\mu_\mathrm{v}}\right)^{\frac{1}{2}}
  \cdot\left(\frac{M_\mathrm{v}}{M_\mathrm{a}}\right)^{\frac{1}{4}}
  \right]^{2}
\end{equation}
\begin{equation}
\Phi_\mathrm{va} = \frac{\sqrt{2}}{4}
  \cdot\left(1+\frac{M_\mathrm{v}}{M_\mathrm{a}}\right)^{-\frac{1}{2}}
  \cdot\left[1
  + \left(\frac{\mu_\mathrm{v}}{\mu_\mathrm{a}}\right)^{\frac{1}{2}}
  \cdot\left(\frac{M_\mathrm{a}}{M_\mathrm{v}}\right)^{\frac{1}{4}}
  \right]^{2}
\end{equation}

% ------------------------------------------------------------
\subsection{Thermal conductivity}

Reid et al.\ [15] suggests the following expression, which was
originally proposed by Wassiljewa [17], as the basis of the
calculation of the thermal conductivity of the mixture,
\begin{equation}
k_\mathrm{m} = \sum_{i=1}^{n}
  \frac{x_i\cdot k_i}{\sum_{j=1}^{n} x_j \Theta_{ij}}
\end{equation}
Based also on the original investigations by Mason and Saxena [18],
they have recommended that
\begin{equation}
\Theta_{ij} = \epsilon\cdot\frac{\left[1
  + \left(\dfrac{k_{\mathrm{tr}i}}{k_{\mathrm{tr}j}}\right)^{1/2}
  \cdot\left(\dfrac{M_i}{M_j}\right)^{1/4}\right]^{2}}
  {\left[8\left(1+\dfrac{M_i}{M_j}\right)\right]^{1/2}}
\end{equation}
where the ratio of the monoatomic values of thermal conductivity in the
previous expression is calculated according to Ref.~[15] by
\begin{equation}
\frac{k_{\mathrm{tr}i}}{k_{\mathrm{tr}j}}
 = \frac{\mu_i}{\mu_j}\cdot\frac{M_j}{M_i}
\end{equation}

% [Prose -- p.\ 1101.  Mason and Saxena [18] recommend
%  epsilon = 1.065 for non-polar gases; Tondon and Saxena [19]
%  suggest 0.85 for polar/non-polar mixtures; Reid et al. [15]
%  recommend unity, which is adopted here.]

\noindent Under this assumption, the substitution of Eq.~(26) into
Eq.~(25) leads to
\begin{equation}
\Phi_{ij} = \Theta_{ij}
\end{equation}
The following expression, employed for calculation of the thermal
conductivity of mixtures, is derived from Eq.~(24),
\begin{equation}
k_\mathrm{m} = \frac{(1-x_\mathrm{v})\cdot k_\mathrm{a}}
  {(1-x_\mathrm{v}) + x_\mathrm{v}\cdot\Phi_\mathrm{av}}
  + \frac{x_\mathrm{v}\cdot k_\mathrm{v}}
  {x_\mathrm{v} + (1-x_\mathrm{v})\cdot\Phi_\mathrm{va}}
\end{equation}
Taking into consideration Eq.~(5), the previous expression becomes
\begin{equation}
\begin{split}
k_\mathrm{m} = {} &
 \frac{\left[1-\grp\right]\cdot k_\mathrm{a}}
      {\left[1-\grp\right] + \grp\,\Phi_\mathrm{av}} \\[6pt]
 & + \frac{\grp\cdot k_\mathrm{v}}
        {\grp + \left[1-\grp\right]\cdot\Phi_\mathrm{va}}
\end{split}
\end{equation}

% ------------------------------------------------------------
\subsection{Specific heat capacity}

% [Prose -- p.\ 1102, on the linear mixing approach of Wong and
%  Embelton [20] and its adoption by [3], [21], [4] and [6].]

\noindent Following the same approach, the specific heat capacity of
the ideal gas mixture can be expressed as
\begin{equation}
c_{p\mathrm{m}} = c_{p\mathrm{a}}\cdot x_\mathrm{a}
  \cdot\frac{M_\mathrm{a}}{M_\mathrm{m}}
  + c_{p\mathrm{v}}\cdot x_\mathrm{v}
  \cdot\frac{M_\mathrm{v}}{M_\mathrm{m}}
\end{equation}
To account for the real gas behavior, the correction term $\Delta c_p$
was proposed according to Reid et al.\ [15],
\begin{equation}
c_{p,\mathrm{m}} - c_{p,\mathrm{mi}} = \Delta c_p
\end{equation}

% [Prose -- p.\ 1102.  Delta c_p is a complex function of the
%  deviation functions and the acentric factor; per Melling et al.
%  [6] it gives at most a 1.5% correction, so it is neglected and
%  c_pmi = c_pm is used.]

\noindent However, taking into account Eq.~(5), the molar fraction of
dry air is
\begin{equation}
x_\mathrm{a} = 1 - x_\mathrm{v}
             = 1 - \fPT\cdot\RH\cdot\frac{\Psv}{P_0}
\end{equation}
and since
\begin{equation}
M_\mathrm{m} = M_\mathrm{a}\cdot x_\mathrm{a}
             + M_\mathrm{v}\cdot x_\mathrm{v}
\end{equation}
Eq.~(30) becomes
\begin{equation}
c_{p\mathrm{m}} = \frac{
   \begin{array}{c}
   c_{p\mathrm{a}}\cdot\left[1-\grp\right]\cdot M_\mathrm{a} \\
   + \,c_{p\mathrm{v}}\cdot\grp\cdot M_\mathrm{v}
   \end{array}}
  {M_\mathrm{a}\cdot\left[1-\grp\right]
       + M_\mathrm{v}\cdot\grp}
\end{equation}

% ------------------------------------------------------------
\subsection{Thermal diffusivity}

Thermal diffusivity is calculated from its definition expression,
\begin{equation}
\alpha_\mathrm{m} = \frac{k_\mathrm{m}}
  {\rho_\mathrm{m}\cdot c_{p\mathrm{m}}}
\end{equation}
taking into account the previously derived Eqs.~(12), (29) and (34) for
the respective calculations of density, thermal conductivity and
specific heat capacity of humid air.

% ------------------------------------------------------------
\subsection{Prandtl number}

This dimensionless number is defined as
$Pr_\mathrm{m} = \nu_\mathrm{m}/\alpha_\mathrm{m}$. Since the thermal
diffusivity from Eq.~(35) and the kinematic viscosity, as derived from
$\mu_\mathrm{m} = \rho_\mathrm{m}\cdot\nu_\mathrm{m}$, are both
functions of density, the Prandtl number is evaluated as a function of
known thermophysical properties $\mu_\mathrm{m}$, $c_{p\mathrm{m}}$
and $k_\mathrm{m}$, as
\begin{equation}
Pr_\mathrm{m} = \frac{\mu_\mathrm{m}\cdot c_{p,\mathrm{m}}}
                     {k_\mathrm{m}}
\end{equation}
with the values of properties $\mu_\mathrm{m}$, $c_{p,\mathrm{m}}$ and
$k_\mathrm{m}$ derived from Eqs.~(21), (34) and (29), respectively.

% ============================================================
\section{The thermophysical and transport properties of dry air and
         water vapor}
% ============================================================

% [Prose -- p.\ 1103.]

\noindent As soon as the molar mass $M_{a,i}$ and the molar fraction
$x_{a,i}$ of the $n$ individual constituent gases of dry air mixture
are known, the molar mass of atmospheric air can be derived by the
following expression,
\begin{equation}
M_\mathrm{a} = \frac{\sum_{i=1}^{n} x_{a,i}\cdot M_{a,i}}
                    {\sum_{i=1}^{n} x_{a,i}}
\end{equation}
Assuming a composition similar to that of standard dry air, as
described for example by Ref.~[22], a molar mass of
$M_\mathrm{a} = 28.9635$ kg/kmol is derived from the previous
expression.

% [Prose -- p.\ 1103.  All properties were taken from the Handbook
%  of Heat Transfer [23], as compiled from Irvine and Liley [24].]

\noindent The viscosity of dry air in Ns/m$^2\times10^{-6}$ is offered
by the following correlation,
\begin{equation}
\mu_\mathrm{a} = \mathrm{MA}_0 + \mathrm{MA}_1\cdot T
  + \mathrm{MA}_2\cdot T^{2} + \mathrm{MA}_3\cdot T^{3}
  + \mathrm{MA}_4\cdot T^{4}
\end{equation}
in the temperature range $-23\,\mdegC \leqslant t \leqslant 327\,\mdegC$,
for the following values of the numerical constants
$\mathrm{MA}_0 = -9.8601\times10^{-1}$,
$\mathrm{MA}_1 = 9.080125\times10^{-2}$,
$\mathrm{MA}_2 = -1.17635575\times10^{-4}$,
$\mathrm{MA}_3 = 1.2349703\times10^{-7}$ and
$\mathrm{MA}_4 = -5.7971299\times10^{-11}$.

\noindent The thermal conductivity of dry air in
W/m\,K\,$\times10^{-3}$ is expressed by the following correlation,
% CAUTION: the "x 10^-3" in the line above is as printed, but these
% coefficients already return W/m K directly -- Eq. (39) evaluates
% to 0.02403 at T = 273.15 K, the correct dry-air value.  Do NOT
% apply a further 1e-3 factor here.  Eq. (42) for water vapor DOES
% require it.
\begin{equation}
k_\mathrm{a} = \mathrm{KA}_0 + \mathrm{KA}_1\cdot T
  + \mathrm{KA}_2\cdot T^{2} + \mathrm{KA}_3\cdot T^{3}
  + \mathrm{KA}_4\cdot T^{4} + \mathrm{KA}_5\cdot T^{5}
\end{equation}
for
$\mathrm{KA}_0 = -2.276501\times10^{-3}$,
$\mathrm{KA}_1 = 1.2598485\times10^{-4}$,
$\mathrm{KA}_2 = -1.4815235\times10^{-7}$,
$\mathrm{KA}_3 = 1.73550646\times10^{-10}$,
$\mathrm{KA}_4 = -1.066657\times10^{-13}$ and
$\mathrm{KA}_5 = 2.47663035\times10^{-17}$.
This expression is valid at the temperature range between
$-23\,\mdegC \leqslant t \leqslant 777\,\mdegC$.

\noindent The specific heat capacity of dry air in kJ/kg\,K, which is
valid for the same temperature range, is given by the following
expression
\begin{equation}
c_{p\mathrm{a}} = \mathrm{CA}_0 + \mathrm{CA}_1\cdot T
  + \mathrm{CA}_2\cdot T^{2} + \mathrm{CA}_3\cdot T^{3}
  + \mathrm{CA}_4\cdot T^{4}
\end{equation}
for the following values of the numerical constants
$\mathrm{CA}_0 = 0.103409\times10^{1}$,
$\mathrm{CA}_1 = -0.284887\times10^{-3}$,
$\mathrm{CA}_2 = 0.7816818\times10^{-6}$,
$\mathrm{CA}_3 = -0.4970786\times10^{-9}$ and
$\mathrm{CA}_4 = 0.1077024\times10^{-12}$.

% [Prose -- p.\ 1103.  Water vapor data from the same source as
%  compiled from Toulukian et al. [25], fitted for
%  0 <= t <= 120 degC.]

\noindent The viscosity in Ns/m$^2\times10^{-6}$ was determined by the
following linear expression,
\begin{equation}
\mu_\mathrm{v} = \mathrm{MV}_0 + \mathrm{MV}_1\cdot t
\end{equation}
with $\mathrm{MV}_0 = 8.058131868\times10^{1}$ and
$\mathrm{MV}_1 = 4.000549451\times10^{-1}$.

\noindent The thermal conductivity in W/m\,K\,$\times10^{-3}$ was
determined by the expression,
\begin{equation}
k_\mathrm{v} = \mathrm{KV}_0 + \mathrm{KV}_1 t
  + \mathrm{KV}_2\cdot t^{2}
\end{equation}
where $\mathrm{KV}_0 = 1.761758242\times10^{1}$,
$\mathrm{KV}_1 = 5.558941059\times10^{-2}$ and
$\mathrm{KV}_2 = -1.663336663\times10^{-4}$.

\noindent The specific heat capacity in (kJ/kg\,K) was determined by
the following expression,
\begin{equation}
c_{p\mathrm{v}} = \mathrm{CV}_0 + \mathrm{CV}_1\cdot t
  + \mathrm{CV}_2\cdot t^{2}
\end{equation}
with $\mathrm{CV}_0 = 1.86910989$,
$\mathrm{CV}_1 = -2.578421578\times10^{-4}$ and
$\mathrm{CV}_2 = 1.941058941\times10^{-5}$.

% ============================================================
\section{Results and discussion}
% ============================================================
% [Prose -- pp.\ 1103-1105.  Quantitative results:
%  DENSITY (Fig. 1): RH 0->100% gives negligible change near
%   freezing, about 4.8% decrease at 50 degC, about 37.5% at 100 degC.
%  VISCOSITY (Fig. 2): about 6.8% decrease at 50 degC, about 45% at
%   100 degC; each RH curve has a viscosity maximum that moves to
%   higher temperature as RH falls -- about 86 degC for the
%   2e-5 Ns/m2 maximum at RH = 20%, about 40 degC for the
%   1.8e-5 Ns/m2 maximum at RH = 100%.
%  THERMAL CONDUCTIVITY (Fig. 3): 3.5% decrease at 50 degC, about
%   21.5% at 100 degC; the maximum moves from 63 degC to about
%   94 degC as RH falls from saturation to about 40%.
%  SPECIFIC HEAT CAPACITY (Fig. 4): increase of about 7% at
%   50 degC and 100% at 100 degC.
%  Figs. 5 and 6 show thermal diffusivity and Prandtl number.]

\noindent The density of the saturated mixture for the temperature
range between 0 and 100 \degC\ was fitted by the following third
degree polynomial,
\begin{equation}
\rho_\mathrm{m} = \mathrm{SD}_0 + \mathrm{SD}_1\cdot t
  + \mathrm{SD}_2\cdot t^{2} + \mathrm{SD}_3\cdot t^{3}
\end{equation}

\noindent The saturated mixture viscosity for the temperature range
between 0 and 100 \degC\ was fitted by the following fourth degree
polynomial,
\begin{equation}
\mu_\mathrm{m} = \mathrm{SV}_0 + \mathrm{SV}_1\cdot t
  + \mathrm{SV}_2\cdot t^{2} + \mathrm{SV}_3\cdot t^{3}
  + \mathrm{SV}_4\cdot t^{4}
\end{equation}

\noindent The thermal conductivity of the saturated mixture for the
temperature range of interest was fitted by the following fourth
degree polynomial,
\begin{equation}
k_\mathrm{m} = \mathrm{SK}_0 + \mathrm{SK}_1\cdot t
  + \mathrm{SK}_2\cdot t^{2} + \mathrm{SK}_3\cdot t^{3}
  + \mathrm{SK}_4\cdot t^{4}
\end{equation}

\noindent The saturated mixture specific heat capacity for the
temperature range between 0 and 100 \degC\ was fitted by the following
fifth degree polynomial,
\begin{equation}
c_{p\mathrm{m}} = \mathrm{SC}_0 + \mathrm{SC}_1\cdot t
  + \mathrm{SC}_2\cdot t^{2} + \mathrm{SC}_3\cdot t^{3}
  + \mathrm{SC}_4\cdot t^{4} + \mathrm{SC}_5\cdot t^{5}
\end{equation}

\noindent The corresponding saturated mixture properties for the
temperature range of interest were fitted by the following fourth
degree polynomial,
\begin{equation}
\alpha_\mathrm{m} = \mathrm{SA}_0 + \mathrm{SA}_1\cdot t
  + \mathrm{SA}_2\cdot t^{2} + \mathrm{SA}_3\cdot t^{3}
  + \mathrm{SA}_4\cdot t^{4}
\end{equation}

\noindent The Prandtl number for the saturated mixture for the same
temperature range was fitted by the following fourth degree
polynomial,
\begin{equation}
Pr_\mathrm{m} = \mathrm{SP}_0 + \mathrm{SP}_1\cdot t
  + \mathrm{SP}_2\cdot t^{2} + \mathrm{SP}_3\cdot t^{3}
  + \mathrm{SP}_4\cdot t^{4}
\end{equation}
with a coefficient of determination and corresponding numerical
constants as shown in Table~4.

% ---------------- Figures 1-6 ----------------
\begin{figure}[htbp]\centering
% \includegraphics[width=0.6\textwidth]{fig1.pdf}
\framebox[0.6\textwidth][c]{\rule{0pt}{5cm} Fig.~1}
\caption*{Fig.~1. The moist air density as a function of temperature
with the relative humidity as a parameter ranging between dry air
(top curve RH = 0\%) and saturation conditions (lower curve
RH = 100\%) in 10\% steps.}
\end{figure}

\begin{figure}[htbp]\centering
% \includegraphics[width=0.6\textwidth]{fig2.pdf}
\framebox[0.6\textwidth][c]{\rule{0pt}{5cm} Fig.~2}
\caption*{Fig.~2. The viscosity of moist air as a function of
temperature with the relative humidity as a parameter ranging between
dry air (top curve RH = 0\%) and saturation conditions (lower curve
RH = 100\%) in 10\% steps.}
\end{figure}

\begin{figure}[htbp]\centering
% \includegraphics[width=0.6\textwidth]{fig3.pdf}
\framebox[0.6\textwidth][c]{\rule{0pt}{5cm} Fig.~3}
\caption*{Fig.~3. The thermal conductivity of moist air as a function
of temperature with the relative humidity as a parameter ranging
between dry air (top curve RH = 0\%) and saturation conditions (lower
curve RH = 100\%) in 10\% steps.}
\end{figure}

\begin{figure}[htbp]\centering
% \includegraphics[width=0.6\textwidth]{fig4.pdf}
\framebox[0.6\textwidth][c]{\rule{0pt}{5cm} Fig.~4}
\caption*{Fig.~4. The specific heat capacity of moist air as a function
of temperature with the relative humidity as a parameter ranging
between dry air (top curve RH = 0\%) and saturation conditions (lower
curve RH = 100\%) in 10\% steps.}
\end{figure}

\begin{figure}[htbp]\centering
% \includegraphics[width=0.6\textwidth]{fig5.pdf}
\framebox[0.6\textwidth][c]{\rule{0pt}{5cm} Fig.~5}
\caption*{Fig.~5. The thermal diffusivity of moist air as a function of
temperature with the relative humidity as a parameter ranging between
dry air (top curve RH = 0\%) and saturation conditions (lower curve
RH = 100\%) in 10\% steps.}
\end{figure}

\begin{figure}[htbp]\centering
% \includegraphics[width=0.6\textwidth]{fig6.pdf}
\framebox[0.6\textwidth][c]{\rule{0pt}{5cm} Fig.~6}
\caption*{Fig.~6. The Prandtl number of moist air as a function of
temperature with the relative humidity as a parameter ranging between
dry air (lower curve RH = 0\%) and saturation conditions (top curve
RH = 100\%) in 10\% steps.}
\end{figure}

% ============================================================
\section{Comparisons with the results from earlier investigations}
% ============================================================

% [Prose -- p.\ 1106.  Nelson's report [7] is the single readily
%  available source of formulae for this temperature range, but
%  citations and derivations are largely absent from it.]

\noindent The recommended expression for mixture density is
\begin{equation}
\rho_\mathrm{m} = (3.484 - 1.317\cdot x_\mathrm{v})
  \cdot\frac{P_0}{273.15+t}
\end{equation}

% [Prose -- p.\ 1106.  Eq. (50) follows from Eq. (12) with
%  z_m = f(P,T) = 1 and R = 8.314 J/mol K,
%  M_a = 28.963 kg/kmol, M_v = 18.02 kg/kmol, giving identical
%  results to Eq. (12).]

\noindent The proposed expression for mixture viscosity is
\begin{equation}
\mu_\mathrm{m} = \frac{\mu_\mathrm{a}
  + x_\mathrm{v}\cdot(0.7887\cdot\mu_\mathrm{v}-\mu_\mathrm{a})}
  {1 - 0.2113\cdot x_\mathrm{v}}
\end{equation}
for which the report referred to Ref.~[26], where its derivation is
attributed to an earlier cited work by Herning and Zipperer
recommending the derivation of mixture viscosity as
\begin{equation}
\mu_\mathrm{m} = \frac{\sum_{i=1}^{n} x_i\cdot\mu_i\cdot(M_i)^{1/2}}
                      {\sum_{i=1}^{n} x_i\cdot(M_i)^{1/2}}
\end{equation}

% [Prose -- p.\ 1106.  Applied to a two-component system this is
%  identical to Eq. (9) in Ref. [6], as proposed by Krischer and
%  Kast [27], and completely different from Eq. (21).]

\noindent The proposed expression for the specific heat capacity,
although of an unspecified origin is
\begin{equation}
c_{p\mathrm{m}} = c_{p\mathrm{a}}
  + 0.622\cdot\frac{P_\mathrm{v}}{P_0-P_\mathrm{v}}
\end{equation}

% [Prose -- p.\ 1106.  Eq. (53) follows from Eq. (34) by omitting
%  the vapor-pressure term of the denominator and assuming unity
%  water vapor specific heat capacity, which actually ranges
%  between about 1.86 and 2.03 kJ/kg K over the range of interest;
%  hence Eq. (53) is of questionable accuracy.]

\noindent The proposed thermal conductivity expression in (mW/m\,K)
units is
\begin{equation}
k_\mathrm{m} = \frac{k_\mathrm{a}
  + x_\mathrm{v}\cdot(0.8536\cdot k_\mathrm{v}-k_\mathrm{a})}
  {1 - 0.1464\cdot x_\mathrm{v}}
\end{equation}
This is also attributed to the chemical engineers handbook [26] in
which its derivation is attributed to the following expression,
recommended by Friend and Adler [28],
\begin{equation}
k_\mathrm{m} = \frac{\sum_{i=1}^{n} x_i\cdot k_i\cdot(M_i)^{1/3}}
                    {\sum_{i=1}^{n} x_i\cdot(M_i)^{1/3}}
 = \frac{x_\mathrm{a}\cdot k_\mathrm{a}\cdot(M_\mathrm{a})^{1/3}
       + x_\mathrm{v}\cdot k_\mathrm{v}\cdot(M_\mathrm{v})^{1/3}}
       {x_\mathrm{a}\cdot(M_\mathrm{a})^{1/3}
       + x_\mathrm{v}\cdot(M_\mathrm{v})^{1/3}}
\end{equation}

% [Prose -- pp.\ 1106-1108.  Discussion of the measured data:
%  agreement with Nelson's viscosity correlation within about 4%;
%  measurements typically under 8% above the calculated values;
%  thermal conductivity deviations from Nelson about 1%; measured
%  conductivities up to about 7% high at 65 degC and 10% at
%  80 degC.  Nelson's c_p correlation gives unrealistically high
%  values at high RH and temperature.]

% ---------------- Table 2 ----------------
\begin{table}[htbp]
\centering
\footnotesize
\caption*{Table 2\\ Data corresponding to viscosity measurements from
          various literature sources ($\times10^{-6}$ kg/m\,s)}
\setlength{\tabcolsep}{4pt}
\begin{tabular}{llllllllllll}
\toprule
\multicolumn{12}{l}{\textit{From Kestin and Whitelaw} [6]} \\
25 \degC & 18.451 & 18.446 & 18.441 & 18.419 & 18.2 & & & & & & \\
         & 19\%   & 25.5\% & 35\%   & 51\%   & 54\% & & & & & & \\
50 \degC & 19.593 & 19.591 & 19.575 & 19.539 & 19.474 & 19.247 & & & & & \\
         & 15.5\% & 20\%   & 25\%   & 34\%   & 51\%   & 98\%   & & & & & \\
75 \degC & 20.632 & 20.588 & 20.497 & 20.357 & 20.046 & 19.586 & 19.252 & 18.792 & 18.781 & & \\
         & 14\%   & 20\%   & 25\%   & 35\%   & 51\%   & 70\%   & 83\%   & 97\%   & 100\%  & & \\
\midrule
\multicolumn{12}{l}{\textit{From Hochrainer and Munczak} [6]} \\
20 \degC & 18.176 & 18.150 & 18.136 & 18.134 & & & & & & & \\
         & 0\%    & 62\%   & 82\%   & 93\%   & & & & & & & \\
30 \degC & 18.647 & 18.620 & 18.617 & 18.586 & 18.569 & & & & & & \\
         & 0\%    & 41\%   & 61\%   & 81\%   & 92\%   & & & & & & \\
40 \degC & 19.111 & 19.111 & 19.080 & 19.053 & 19.017 & 18.995 & & & & & \\
         & 0\%    & 0\%    & 41\%   & 63\%   & 83\%   & 93\%   & & & & & \\
50 \degC & 19.588 & 19.553 & 19.483 & 19.426 & 19.363 & 19.343 & & & & & \\
         & 0\%    & 21\%   & 41\%   & 61\%   & 82\%   & 93\%   & & & & & \\
\midrule
\multicolumn{12}{l}{\textit{From Vargaftik} [6,29]} \\
50 \degC  & 19.550 & 19.140 & & & & & & & & & \\
          & 0\%    & 82\%   & & & & & & & & & \\
60 \degC  & 20.01  & 19.600 & 19.04 & & & & & & & & \\
          & 0\%    & 51\%   & 100\% & & & & & & & & \\
70 \degC  & 20.46  & 20.05  & 19.500 & 18.800 & & & & & & & \\
          & 0\%    & 32\%   & 65\%   & 97\%   & & & & & & & \\
80 \degC  & 20.910 & 20.510 & 19.950 & 19.250 & 18.430 & 17.500 & & & & & \\
          & 0\%    & 21\%   & 43\%   & 64\%   & 86\%   & 100\%  & & & & & \\
90 \degC  & 21.350 & 20.950 & 20.350 & 19.690 & 18.870 & 17.920 & 16.890 & 15.770 & & & \\
          & 0\%    & 14\%   & 29\%   & 43\%   & 58\%   & 72\%   & 87\%   & 100\%  & & & \\
100 \degC & 21.800 & 21.400 & 20.840 & 20.140 & 19.310 & 18.360 & 17.320 & 16.180 & 14.990 & 13.750 & 12.470 \\
          & 0\%    & 10\%   & 20\%   & 30\%   & 40\%   & 50\%   & 60\%   & 70\%   & 80\%   & 90\%   & 100\% \\
\midrule
\multicolumn{12}{l}{\textit{From Mason and Monchick} [30]} \\
25 \degC & 18.300 & 18.000 & & & & & & & & & \\
         & 0\%    & 100\%  & & & & & & & & & \\
50 \degC & 19.550 & 19.100 & 19.070 & & & & & & & & \\
         & 0\%    & 82\%   & 100\%  & & & & & & & & \\
75 \degC & 20.850 & 20.170 & 19.660 & 19.100 & 18.400 & & & & & & \\
         & 0\%    & 26\%   & 53\%   & 79\%   & 100\%  & & & & & & \\
\bottomrule
\end{tabular}
\end{table}

% ---------------- Table 3 ----------------
\begin{table}[htbp]
\centering
\footnotesize
\caption*{Table 3\\ Data corresponding to thermal conductivity
          measurements from various literature sources
          ($\times10^{-2}$ W/m\,K)}
\setlength{\tabcolsep}{4pt}
\begin{tabular}{lllllllllll}
\toprule
\multicolumn{11}{l}{\textit{From Gruss and Schmick} [31]} \\
80 \degC & 2.9845 & 3.044 & 3.088 & 3.094 & 3.091 & 3.088 & 3.0949 & 3.062 & 3.074 & 2.9810 \\
         & 0\%    & 15\%  & 32\%  & 37\%  & 42\%  & 48\%  & 53\%   & 65\%  & 67\%  & 95\% \\
\midrule
\multicolumn{11}{l}{\textit{From Vargaftik} [29]} \\
80 \degC & 2.989 & 3.069 & 3.103 & 3.078 & 3.019 & & & & & \\
         & 0\%   & 21\%  & 43\%  & 65\%  & 86\%  & & & & & \\
\midrule
\multicolumn{11}{l}{\textit{From Touloukian et al.} [32,6]} \\
60 \degC & 2.92 & 2.96 & 2.96 & 2.92    & & & & & & \\
         & 0\%  & 40\% & 81\% & >100\%  & & & & & & \\
\bottomrule
\end{tabular}
\end{table}

% ---------------- Figures 7-9 ----------------
\begin{figure}[htbp]\centering
% \includegraphics[width=0.6\textwidth]{fig7.pdf}
\framebox[0.6\textwidth][c]{\rule{0pt}{5cm} Fig.~7}
\caption*{Fig.~7. Comparative viscosity of moist air at various
temperature levels according to results from the present analysis
(curves in solid lines), Nelsons correlation (curves in broken lines)
and discrete data from earlier measurements in the literature.}
\end{figure}

\begin{figure}[htbp]\centering
% \includegraphics[width=0.6\textwidth]{fig8.pdf}
\framebox[0.6\textwidth][c]{\rule{0pt}{5cm} Fig.~8}
\caption*{Fig.~8. Comparative thermal conductivity of moist air at
various temperature levels according to results from the present
analysis (curves in solid lines), Nelsons correlation (curves in
broken lines) and discrete data from earlier measurements in the
literature.}
\end{figure}

\begin{figure}[htbp]\centering
% \includegraphics[width=0.6\textwidth]{fig9.pdf}
\framebox[0.6\textwidth][c]{\rule{0pt}{5cm} Fig.~9}
\caption*{Fig.~9. Comparative specific heat capacity of moist air at
various temperature levels according to results from the present
analysis (curves in solid lines) and Nelsons correlation (53) (curves
in broken lines) for a relative humidiry ranging between RH = 0\%
(lower curves) to RH = 100\% (top curves).}
\end{figure}

% ---------------- Table 4 ----------------
% IMPORTANT: SK2 below is printed in the original as
% -1.788037411E-2.  That is a typo.  Substituting it makes
% k_m = -178 W/m K at 100 degC.  The intended value is
% -1.788037411E-7, verified numerically against Eq. (29):
%    t (degC)      0       20       40       60       80      100
%    Eq. (29)  0.02403  0.02547  0.02665  0.02724  0.02617  0.02147
%    poly E-7  0.02401  0.02538  0.02646  0.02701  0.02663  0.02482
% which also matches Fig. 3 and the stated conductivity maximum
% near 63 degC.  The printed value is retained in the table for
% fidelity; use the corrected one in any computation.
\begin{landscape}
\begin{table}[htbp]
\centering
\scriptsize
\caption*{Table 4\\ The numerical constants and coefficients of
determination (COD) for the proposed polynomial fit expressions for the
following saturated mixture properties, density, viscosity, thermal
conductivity, specific heat capacity, thermal diffusivity and Prandtl
number}
\setlength{\tabcolsep}{3pt}
\begin{tabular}{llllll}
\toprule
Density (kg/m$^3$) & Viscosity (Ns/m$^2$)
 & Thermal conductivity (W/m K) & Specific heat capacity (kJ/kg K)
 & Thermal diffusivity (m$^2$/s) & Prandtl number \\
COD = 0.999954 & COD = 0.999997 & COD = 0.999985
 & COD = 0.999905 & COD = 0.999900 & COD = 0.999990 \\
\midrule
$\mathrm{SD}_0 = 1.293393662$
 & $\mathrm{SV}_0 = 1.715747771$E-5
 & $\mathrm{SK}_0 = 2.40073953$E-2
 & $\mathrm{SC}_0 = 1.004571427$
 & $\mathrm{SA}_0 = 1.847185729$E-5
 & $\mathrm{SP}_0 = 0.7215798365$ \\
$\mathrm{SD}_1 = -5.538444326$E-3
 & $\mathrm{SV}_1 = 4.722402075$E-8
 & $\mathrm{SK}_1 = 7.278410162$E-5
 & $\mathrm{SC}_1 = 2.05063275$E-3
 & $\mathrm{SA}_1 = 1.161914598$E-7
 & $\mathrm{SP}_1 = -3.703124976$E-4 \\
$\mathrm{SD}_2 = 3.860201577$E-5
 & $\mathrm{SV}_2 = -3.663027156$E-10
 & $\mathrm{SK}_2 = -1.788037411$E-2
 & $\mathrm{SC}_2 = -1.631537093$E-4
 & $\mathrm{SA}_2 = 2.373056947$E-10
 & $\mathrm{SP}_2 = 2.240599044$E-5 \\
$\mathrm{SD}_3 = -5.2536065$E-7
 & $\mathrm{SV}_3 = 1.873236686$E-12
 & $\mathrm{SK}_3 = -1.351703529$E-9
 & $\mathrm{SC}_3 = 6.2123003$E-6
 & $\mathrm{SA}_3 = -5.769352751$E-12
 & $\mathrm{SP}_3 = -4.162785412$E-7 \\
 & $\mathrm{SV}_4 = -8.050218737$E-14
 & $\mathrm{SK}_4 = -3.322412767$E-11
 & $\mathrm{SC}_4 = -8.830478888$E-8
 & $\mathrm{SA}_4 = -6.369279936$E-14
 & $\mathrm{SP}_4 = 4.969218948$E-9 \\
 & & & $\mathrm{SC}_5 = 5.071307038$E-10 & & \\
\bottomrule
\end{tabular}
\end{table}
\end{landscape}

% ============================================================
\section{Conclusions}
% ============================================================
% [Conclusions prose -- p.\ 1109.  A complete account of moist-air
%  thermophysical and transport properties for 0-100 degC at
%  101.3 kPa; saturated mixture properties fitted for computerized
%  calculation; with the exception of the specific heat capacity
%  there is very good agreement with previous analyses, and good
%  agreement with the scarce earlier measurements.]

% ============================================================
\begin{thebibliography}{99}
\small

\bibitem{1} Giacomo P. Equation for the determination of density of
moist air (1981). Metrologia 1982;18:33--40.

\bibitem{2} Davies RS. Equation for the determination of density of
moist air (1981/1991). Metrologia 1992;29:67--70.

\bibitem{3} Zuckerwar AJ, Meredith RW. Low-frequency absorption of
sound in air. J Acoust Soc Am 1985;78(3).

\bibitem{4} Rasmussen K. Calculation methods for the physical
properties of air used in the calibration of microphones. Report
PL-11b. Department of Acoustic Technology, Techn. University of
Denmark; 1997.

\bibitem{5} Hyland RW, Wexler A. Formulations for the thermodynamic
properties of dry air from 173.15 K to 473.15 K and of saturated moist
air from 173.15 to 372.15 K at pressures to 5 MPa. Trans ASHRAE
1983;89(IIa):520--35.

\bibitem{6} Melling A, Noppenberger S, Still M, Venzke H.
Interpolation correlations for fluid properties of humid air in the
temperature range 100 \degC\ to 200 \degC. J Phys Chem Ref Data
1997;26(4):1111--23.

\bibitem{7} Nelson R. Material properties in SI units, part 4. Chem
Eng Prog 1980;76:83.

\bibitem{8} Hyland RW, Wexler A. Formulations for the thermodynamic
properties of the saturated phases of H$_2$O from 173.15 K to
473.15 K. Trans ASHRAE 1983;89(IIa):520--35.

\bibitem{9} Hyland RW, Wexler A. The second interaction (cross) virial
coefficient for moist air. J Res NBS 1973;77A(1):133--47.

\bibitem{10} Hyland RW. A correllation for the second interaction
virial coefficients and enhancement factors for CO$_2$-free moist air
from $-50$ to 90 \degC. J Res National Bureau Standards A Phys Chem
1975;79A(4):551--60.

\bibitem{11} Hardy B. ITS-90 formulations for vapor pressure, frost
point temperature, dew point temperature and enhancement factors in
the range $-100$ to 100 \degC. Proc 3rd Int Sympos Humidity Moisture.
London: National Physical Lab (N.P.L.); 1998. p.\ 214--21.

\bibitem{12} Greenspan L. Functional equations for the enhancement
factors for CO$_2$--free moist air. J Res National Bureau
Standards---A Phys Chem 1976;80A(1):41--4.

\bibitem{13} Alduchov OA, Eskridge RE. Improved magnus form
approximation of saturation vapor pressure. Technical report number
DOE/ER/61011-T6/Nov. USDOE Office of Energy Research: Washington, DC;
1997.

\bibitem{14} Van Wylen GJ, Sonntag RE. Fundamentals of classical
thermodynamics. NY: Wiley; 1973.

\bibitem{15} Reid RC, Prausnitz JM, Poling BE. The properties of gases
and liquids, Chemical Engineering Series. McGraw Hill Int. Editions;
1988.

\bibitem{16} Wilke CR. A viscosity equation for gas mixtures. J Chem
Phys 1950;18:517.

\bibitem{17} Wassiljewa A. Physik Zeitshrift 1904;5:734.

\bibitem{18} Mason EA, Saxena SC. Approximate formula for the thermal
conductivity of gas mixtures. Phys Fluids 1958;1(5):361--9.

\bibitem{19} Tondon PK, Saxena SC. Calculation of thermal conductivity
of polar--nonpolar gas mixtures. Appl Sci Res 1968;19:163--70.

\bibitem{20} Wong GSK, Embleton TFW. Variation of specific heats and
of specific heat ratio in air with humidity. J Acoust Soc Am
1984;76(2):555--9.

\bibitem{21} Durst F, Noppenberger S, Still M, Venzke H. Influence of
humidity on hot wire measurements. Meas Sci Technol 1996;7:1517--28.

\bibitem{22} U.S. Standard Atmosphere 1976. U.S. Government Printing
Office: Washington, DC; 1976. p.\ 3--33.

\bibitem{23} Rohsenow WM, Hartnett JP, Cho YI. Handbook of heat
transfer. 3rd ed. McGraw-Hill; 1998.

\bibitem{24} Irvine TF, Liley P. Steam and gas tables with computer
equations. San Diego: Academic Press; 1984.

\bibitem{25} Touloukian YS, Powell RW, Ho CY, Clemend PG.
Thermophysical properties of matter, vol.\ 1. NY; 1970.

\bibitem{26} Perry RH, Chilton CH. Chemical engineers handbook. 5th
ed. NY: McGraw-Hill; 1973.

\bibitem{27} Krischer O, Kast W. 3rd ed. Drying technology-the
scientific background of drying technology, vol.\ 1. Berlin,
Heidelberg: Springer; 1978.

\bibitem{28} Fried C, Adler J. In: Cambel, Fenn, editors. Transport
properties in gases. Evanston, Ill: Northwestern University Press;
1958. p.\ 128--31.

\bibitem{29} Vargaftik NB. Handbook of physical properties of liquids
and gases-pure substances and mixtures. Washington: Hemisphere; 1983.

\bibitem{30} Kestin J, Whitelaw JH. The viscosity of dry and humid
air. Int J Heat Mass Transfer 1964:1245--55.

\bibitem{31} Gruss H, Schmick H. Scientific publication of the Siemens
Company, vol.\ 7(202); 1928.

\bibitem{32} Touloukian YS, Liley PE, Saxena SC. Thermal conductivity
non metallic liquids and gases, vol.\ 3. NY: IFI/Plenum Data
Corporation; 1970.

\end{thebibliography}

\end{document}
